 
When training \MethodName{} to handle subgoal images, we need the model to accommodate goals with different delays and different levels of image quality, including images generated by our world model. This requires carefully selecting which subgoals are provided as context to the model when training. We train on a combination of real images from future timesteps of the training trajectory and generated images. We found the following sampling scheme to be effective for selecting the timesteps for the real images:
with probability $0.25$, we sample the end-of-segment images (consistent with the prediction target for the world model), and with probability $0.75$ we sample future images uniformly from 0--4 seconds ahead of the current timestep. In addition to these real images, we mitigate the train-test mismatch between real and generated images by also sampling a large number of subgoal images from the world model, and constructing additional training examples with these \emph{generated} images added into the context of \MethodName{} instead of the real future images.








\section{Prompting \texorpdfstring{\MethodName{}}{pi0.7} at runtime}
\label{sec:runtime}

At runtime, we configure \MethodName{} to run with different forms of context depending on the desired behavior \textbf{without any task-specific post-training}. For any task we always prompt the model with the control mode and episode metadata. 
For choosing the episode metadata, we follow
\begin{itemize}[leftmargin=10pt]
    \item Overall speed: set per-task to the $15^\text{th}$ percentile of the episode length from the task.
    \item Overall quality: always set to 5, which is the highest score.
    \item Mistake: always set to false, meaning no mistake.
\end{itemize}
The subtask instruction $\rawtext_t$ is provided either by a learned high-level language policy or by a human supervisor for coaching (see Sec.~\ref{subsec:hl}).
When the subgoal images are used, we refresh the subgoal images whenever the semantic intent changes (i.e., new $\rawtext_t$), or after $\Delta=4$ seconds have elapsed since the last subgoal image was produced, whichever happens first. See Algorithm~\ref{alg:testtime} for the full workflow. We apply asynchronous inference: the visual subgoal and subtask instruction generation happens in separate threads and the VLA inference always uses the latest ones available.

For all experiments we use 5 denoising steps to generate the 50-step action chunks and execute $\hat H \in \{15, 25\}$ steps out of the chunk. Since each prompt component is trained with dropout, \MethodName{} can also be used with classifier-free guidance (CFG) \cite{ho2022cfg} for any part of the prompt, for example to guide the generated actions toward higher speeds. Concretely, each action denoising step follows
\begin{equation*}
\begin{aligned}
    \nabla_\ba &\log \pi_\theta(\ba_{t:t+H} \vert \bo_t, \mathcal{C}_t) + \\ 
    \beta &( \nabla_\ba \log \pi_\theta(\ba_{t:t+H} \vert \bo_t, \mathcal{C}_t) - \nabla_\ba \log \pi_\theta(\ba_{t:t+H} \vert \bo_t, \mathcal{C}_t^\text{uncond})),
\end{aligned}
\end{equation*}
where $\mathcal{C}_t^\text{uncond}$ denotes the set of context used in ``unconditional'' mode and $\beta$ is the CFG weight. While any part of the context could be dropped out, we apply CFG on the episode metadata to elicit strong performance in dexterous tasks. We use moderate values of $\beta \in \{1.3, 1.7, 2.2\}$.

\begin{algorithm}[t]
\caption{Prompting \MethodName{} at test time}
\label{alg:testtime}
\begin{algorithmic}[1]
\State \textbf{Input:} initial observation $\bo_0$, task instruction $\lang$, episode metadata $m$, control mode $c$
\State Initialize subtask $\rawtext$ (from high-level policy or coaching)
\State $\bg^{\star} \sim p_{\psi}(\bg^{\star} \mid \bo_0, \rawtext, m)$
\State $\mathcal{C} = \{\lang, \rawtext, \bg^{\star}, m, c\}$
\State $\ba_{t:t+H} \sim \pi_{\theta}(\ba \mid \bo_{t-T:t}, \mathcal{C})$ \Comment{Optional: CFG}
\For{$t=0,1,2,\ldots$}
    \If {$\rawtext$ changed \textbf{or} $\Delta$-second timer elapsed}
        \State $\bg^{\star} \sim p_{\psi}(\bg^{\star} \mid \bo_t, \rawtext, m)$ \Comment{Non-blocking (async)}
        \State $\mathcal{C} = \{\lang, \rawtext, \bg^{\star}, m, c\}$
    \EndIf
    \If {$\hat H$ steps elapsed since last inference}
        \State $\ba_{t:t+H} \sim \pi_{\theta}(\ba \mid \bo_{t-T:t}, \mathcal{C},\ba_{t:})$ \Comment{Async w/ RTC}
    \EndIf
    \State Execute $\ba_t$
\EndFor
\end{algorithmic}
\end{algorithm}

